\documentclass[UTF8,a4paper,12pt]{article}
\usepackage{ctex} % 支持中文
\usepackage{amsmath}
\usepackage{amssymb}
\usepackage{geometry}
\usepackage{graphicx}
\usepackage{color}
\usepackage{tcolorbox}

\geometry{left=2.5cm,right=2.5cm,top=2.5cm,bottom=2.5cm}

\title{2025高考全程复习独家课程：圆锥曲线核心知识体系}
\author{第6讲 核心全面梳理笔记}
\date{}

\begin{document}
	
\section{射影几何与极点极线理论}

\section{圆锥曲线六大定义}

\subsection{}

\section{焦半径公式与焦点弦定理}

\section{椭圆双曲线焦点三角形命题点}

\section{三一线模型}

\section{抛物线25条结论推导}

\section{圆锥曲线光学性质与工程技术应用}

\section{点乘双根}

\textbf{适用类型：} 类似 \( x_1, x_2, y_1, y_2 \)，\( (x_1 + t)(x_2 + t) \)，\( (y_1 + t)(y_2 + t) \) 或 \( MA \cdot MB \)，\( |MA| \cdot |MB| \)（其中 \( x_1, x_2, y_1, y_2 \) 是直线与曲线的两个交点的横纵坐标，\( A, B \) 直线与曲线的两个交点）以及可转化为上述结构的问题。

\textbf{理论基础：} 二次函数的双根式，若一元二次方程 \( ax^2 + bx + c = 0 \) (\( a \neq 0 \)) 的两根为 \( x_1, x_2 \)，则
\[
ax^2 + bx + c = a(x - x_1)(x - x_2)
\]

\textbf{具体步骤：} 化双根式 $\rightarrow$ 赋值 $\rightarrow$ 变形代入

\noindent 1. 椭圆 \( C: \frac{x^2}{4} + \frac{y^2}{3} = 1 \)，左顶点为 \( A \)，过左焦点 \( P \) 的直线交椭圆于 \( P, Q \) 两点，求证：直线 \( AP \) 和 \( AQ \) 的斜率乘积是定值。
	

\text{题目：}~椭圆
C:\ \frac{x^{2}}{4}+\frac{y^{2}}{3}=1,
\text{左顶点为 }A,\text{ 过左焦点 }F\text{ 的直线交椭圆于 }P,Q\text{ 两点，求证：直线 }AP, AQ
\text{ 的斜率乘积为定值。}

\text{解：}~由椭圆方程可得
a=2,\ b=\sqrt3,\ c=\sqrt{a^{2}-b^{2}}=1,
\text{故左顶点 }A(-2,0),\text{ 左焦点 }F(-1,0).

设过焦点 $F$ 的一条动直线斜率为 $t\ (\neq0)$，则其方程为
\[
y=t(x+1). \tag{1}
\]
该直线与椭圆相交于 $P(x_1,y_1),Q(x_2,y_2)$ 两点，将 (1) 代入椭圆方程
\[
\frac{x^{2}}{4}+\frac{t^{2}(x+1)^{2}}{3}=1
\]
化简得
\[
\frac{1}{12}\Big[(4t^{2}+3)x^{2}+8t^{2}x+(4t^{2}-12)\Big]=0,
\]
根不受常数因子影响，于是 $x_1,x_2$ 为二次方程
\[
(4t^{2}+3)x^{2}+8t^{2}x+(4t^{2}-12)=0 \tag{2}
\]
的两根。

对任一点 $M(x,y)$（为 $P$ 或 $Q$），由 (1) 得 $y=t(x+1)$，直线 $AM$ 的斜率为
\[
k=\frac{y-0}{x-(-2)}=\frac{t(x+1)}{x+2}.
\]
因此
\[
k_{AP}=\frac{t(x_1+1)}{x_1+2},\qquad
k_{AQ}=\frac{t(x_2+1)}{x_2+2},
\]
从而
\[
k_{AP}\cdot k_{AQ}
=t^{2}\frac{(x_1+1)(x_2+1)}{(x_1+2)(x_2+2)}. \tag{3}
\]

下面用“点乘双根法”计算分子、分母。

由 (2) 可写成
\[
f(x)=(4t^{2}+3)x^{2}+8t^{2}x+(4t^{2}-12)
=(4t^{2}+3)(x-x_1)(x-x_2).
\]
令 $x=-1$，得到
\[
f(-1)=(4t^{2}+3)(-1-x_1)(-1-x_2)=(4t^{2}+3)(x_1+1)(x_2+1),
\]
即
\[
(x_1+1)(x_2+1)=\frac{f(-1)}{4t^{2}+3}.
\]
同理令 $x=-2$，
\[
f(-2)=(4t^{2}+3)(-2-x_1)(-2-x_2)=(4t^{2}+3)(x_1+2)(x_2+2),
\]
故
\[
(x_1+2)(x_2+2)=\frac{f(-2)}{4t^{2}+3}.
\]
这正是“点乘双根法”的应用：通过在二次式两端代入同一点，得到关于两根的乘积型表达式。

直接计算：
\[
\begin{aligned}
	f(-1)
	&=(4t^{2}+3)\cdot1+8t^{2}\cdot(-1)+(4t^{2}-12)  \\
	&=4t^{2}+3-8t^{2}+4t^{2}-12
	=-9, \\[4pt]
	f(-2)
	&=(4t^{2}+3)\cdot4+8t^{2}\cdot(-2)+(4t^{2}-12) \\
	&=16t^{2}+12-16t^{2}+4t^{2}-12
	=4t^{2}.
\end{aligned}
\]
于是
\[
(x_1+1)(x_2+1)=\frac{-9}{4t^{2}+3},\qquad
(x_1+2)(x_2+2)=\frac{4t^{2}}{4t^{2}+3}.
\]

代回 (3) 得
\[
k_{AP}\cdot k_{AQ}
=t^{2}\frac{\dfrac{-9}{4t^{2}+3}}{\dfrac{4t^{2}}{4t^{2}+3}}
=t^{2}\cdot\frac{-9}{4t^{2}}
=-\frac{9}{4}.
\]

可见直线 $AP$ 与 $AQ$ 的斜率乘积为
\[
k_{AP}\cdot k_{AQ}=-\frac{9}{4},
\]
与动直线的斜率 $t$ 无关，为一常数。证明完毕。

	\subsection{平移齐次化法}
	
	
	\section
	
	\section{非对称韦达定理}
	
	
	
\end{document}
